% GENERATED FILE. Edit canonical lessons, then run npm run generate:books.
% canonical-source-hash: fnv1a64:13b7e32a39214288
% canonical-lessons: GE-C189-uebersetzen, GE-C189-umarmen, GE-C189-verabreden, GE-C189-verbieten, GE-C189-verbinden

\chapter{To translate, To hug, To arrange, To forbid, To connect}
\label{ch:ge-to-translate-to-hug-to-arrange-to-forbid-to-connect}
\hlchaptermodality{189}

\begin{chapteropening}
\textbf{By the end of this chapter:} \emph{I can say to translate, to hug, to arrange, to forbid and to connect.}

Complete the last lesson of chapter 189: I can say to translate, to hug, to arrange, to forbid and to connect.
\end{chapteropening}

\section[übersetzen]{übersetzen — to translate}

\label{lesson:GE-C189-uebersetzen}



\begin{quote}
Before the new one: say the German for to separate, then the German for to cross.
\end{quote}

\subsection*{You'll want to know: übersetzen}
\textbf{übersetzen} — \textquotedblleft{}to translate\textquotedblright{}.

\begin{cousinweb}
\textbf{übersetzen} is the word for \textquotedblleft{}over\textquotedblright{} + the verb \textquotedblleft{}to set\textquotedblright{}, Old High German \emph{sezzen}, English \emph{set} (EtymDB 2.1): to set words over into another language, a copy of Latin \emph{trādūcere}.
\end{cousinweb}

\begin{grammarlens}[title={What you've built}]
One more word for this chapter.
\end{grammarlens}

\subsection*{Your turn}
\begin{itemize}
\raggedright
  \item \emph{Say it:} \emph{übersetzen}
  \item \emph{Say it:} \emph{übersetzen}, once more
  \item \emph{Say it:} say \emph{übersetzen}
  \item \emph{Recall:} say the German for to separate, then the German for to cross, then say \emph{übersetzen} again
\end{itemize}

\begin{tcolorbox}[breakable,colback=teal!4,colframe=teal!35!black,title={Before you move on}]
What does übersetzen mean? (\textquotedblleft{}To translate\textquotedblright{}.) Say it once more.
\end{tcolorbox}
\section[umarmen]{umarmen — to hug}

\label{lesson:GE-C189-umarmen}



\begin{quote}
Before the new one: say the German for to cross, then the German for to translate.
\end{quote}

\subsection*{You'll want to know: umarmen}
\textbf{umarmen} — \textquotedblleft{}to hug\textquotedblright{}.

\begin{cousinweb}
\textbf{umarmen} is \emph{um-}, \textquotedblleft{}around\textquotedblright{}, + \emph{Arm} + a verb ending: to put your arms around someone. \emph{Arm} $\leftarrow$ Proto-Germanic \emph{*armaz}, English \emph{arm}. (EtymDB 2.1)
\end{cousinweb}

\begin{grammarlens}[title={What you've built}]
One more word for this chapter.
\end{grammarlens}

\subsection*{Your turn}
\begin{itemize}
\raggedright
  \item \emph{Say it:} \emph{umarmen}
  \item \emph{Say it:} \emph{umarmen}, once more
  \item \emph{Say it:} say \emph{umarmen}
  \item \emph{Recall:} say the German for to cross, then the German for to translate, then say \emph{umarmen} again
\end{itemize}

\begin{tcolorbox}[breakable,colback=teal!4,colframe=teal!35!black,title={Before you move on}]
What does umarmen mean? (\textquotedblleft{}To hug\textquotedblright{}.) Say it once more.
\end{tcolorbox}
\section[verabreden]{verabreden — to arrange (sich verabreden: to make a date)}

\label{lesson:GE-C189-verabreden}



\begin{quote}
Before the new one: say the German for to translate, then the German for to hug.
\end{quote}

\subsection*{You'll want to know: verabreden}
\textbf{verabreden} — \textquotedblleft{}to arrange (sich verabreden: to make a date)\textquotedblright{}.

\begin{cousinweb}
\textbf{verabreden} is \emph{ver-} + \emph{ab-} + the verb \textquotedblleft{}to speak\textquotedblright{}, Old High German \emph{redōn} (EtymDB 2.1): to settle something by talking it over.
\end{cousinweb}

\begin{grammarlens}[title={What you've built}]
One more word for this chapter.
\end{grammarlens}

\subsection*{Your turn}
\begin{itemize}
\raggedright
  \item \emph{Say it:} \emph{verabreden}
  \item \emph{Say it:} \emph{verabreden}, once more
  \item \emph{Say it:} say \emph{verabreden}
  \item \emph{Recall:} say the German for to translate, then the German for to hug, then say \emph{verabreden} again
\end{itemize}

\begin{tcolorbox}[breakable,colback=teal!4,colframe=teal!35!black,title={Before you move on}]
What does verabreden mean? (\textquotedblleft{}To arrange (sich verabreden: to make a date)\textquotedblright{}.) Say it once more.
\end{tcolorbox}
\section[verbieten]{verbieten — to forbid}

\label{lesson:GE-C189-verbieten}



\begin{quote}
Before the new one: say the German for to hug, then the German for to arrange.
\end{quote}

\subsection*{You'll want to know: verbieten}
\textbf{verbieten} — \textquotedblleft{}to forbid\textquotedblright{}.

\begin{cousinweb}
\textbf{verbieten} is \emph{ver-} + an old verb meaning \textquotedblleft{}to bid, to command\textquotedblright{} — the same build as English \emph{forbid}, which EtymDB 2.1 lists beside it with Gothic \emph{faurbiudan}.
\end{cousinweb}

\begin{grammarlens}[title={What you've built}]
One more word for this chapter.
\end{grammarlens}

\subsection*{Your turn}
\begin{itemize}
\raggedright
  \item \emph{Say it:} \emph{verbieten}
  \item \emph{Say it:} \emph{verbieten}, once more
  \item \emph{Say it:} say \emph{verbieten}
  \item \emph{Recall:} say the German for to hug, then the German for to arrange, then say \emph{verbieten} again
\end{itemize}

\begin{tcolorbox}[breakable,colback=teal!4,colframe=teal!35!black,title={Before you move on}]
What does verbieten mean? (\textquotedblleft{}To forbid\textquotedblright{}.) Say it once more.
\end{tcolorbox}
\section[verbinden]{verbinden — to connect, to bandage}

\label{lesson:GE-C189-verbinden}



\begin{quote}
Before the new one: say the German for to arrange, then the German for to forbid.
\end{quote}

\subsection*{You'll want to know: verbinden}
\textbf{verbinden} — \textquotedblleft{}to connect, to bandage\textquotedblright{}.

\begin{cousinweb}
\textbf{verbinden} is \emph{ver-} + the verb \textquotedblleft{}to bind\textquotedblright{}, Old High German \emph{bintan} (EtymDB 2.1), the same verb as English \emph{bind}: to tie together, or to bind a wound.
\end{cousinweb}

\begin{grammarlens}[title={What you've built}]
One more word for this chapter.
\end{grammarlens}

\subsection*{Your turn}
\begin{itemize}
\raggedright
  \item \emph{Say it:} \emph{verbinden}
  \item \emph{Say it:} \emph{verbinden}, once more
  \item \emph{Say it:} say \emph{verbinden}
  \item \emph{Recall:} say the German for to arrange, then the German for to forbid, then say \emph{verbinden} again
\end{itemize}

\begin{tcolorbox}[breakable,colback=teal!4,colframe=teal!35!black,title={Before you move on}]
What does verbinden mean? (\textquotedblleft{}To connect, to bandage\textquotedblright{}.) Say it once more.
\end{tcolorbox}
